\section{E2/E2b: local-state memory and noise coupling}
\label{sec:e2}
\subsection{Question and intervention design}
Immediate noise moments hold the local user matrix fixed. During training,
however, users adapt to the shared item state. A shared perturbation can
therefore change local vectors, which later influence the shared model again.
E2 asks whether that local-state pathway has measurable prediction consequences
after the shared state is restored. It is a mechanism probe, not a proposed
private denoising algorithm.

The study starts from 30 E1 round-50 checkpoints: Full, FixedB-r8, and Two-r8,
at targets one and two, across five seeds. Each receives two independently
seeded Gaussian pulses at amplitudes $\alpha\in\{0.25,1\}$ times its own
per-round coordinate noise SD. This gives 120 probes. The pulse is a fresh
perturbation of a saved state, not the replay or removal of an actual earlier
training draw. One ordinary exposure round allows selected local vectors
to respond before reset branches are constructed.

The unpulsed reference, retained-pulse branch, shared-reset branch, and
local-reset branch then continue for ten rounds. A shared reset copies the
reference shared state while retaining the pulsed branch's local vectors.
A local reset copies the reference local vectors while retaining perturbed
shared parameters. A both-reset construction is an identity control.
Branches share future client, negative-sampling, and coordinate-noise streams.
Observations occur immediately after the pulse, after exposure, and at
post-reset lags 0, 1, 4, and 10. All declared probes are retained.

\subsection{Score decomposition and measurement}
For reference $(P,Q)$ and branch $(P+\Delta P,Q+\Delta Q)$,
\begin{equation}
 \Delta S=P\Delta Q^T+\Delta P Q^T+\Delta P\Delta Q^T.
 \label{eq:decompose}
\end{equation}
At a shared reset, $\Delta Q=0$, so the remaining score difference is exactly
$\Delta P Q^T$. Later differences may involve both states again. Squared
norms of the three terms do not add to total squared norm without cross
inner products; they are not additive causal percentages. The implementation
records all terms and checks numerical closure in float64.

Measurements include score RMS across users and items, effective-matrix
distance, top-10 disagreement on training-unseen candidates, and margin
sign changes on 256 seeded random item pairs. The 942-user diagnostics do
not use validation or test labels. ``Exposed'' users are those selected
in the single exposure round; their identity is held fixed when reporting
later outcomes. The normalized score ratio divides a probe's delayed RMS
by its own initial pulse RMS and is accompanied by the absolute value.

\subsection{Initial findings and why a control was needed}
At full pulse amplitude and lag ten after shared reset, Full has mean
top-10 disagreement of $0.5212$\% at target one and $0.3206$\% at target
two. FixedB-r8 has $0.3747$\% and $0.2038$\%. These modest effects show
that restoring the shared item state does not automatically restore every
local prediction. They do not establish an important relevance loss.

Raw-coordinate Two-r8 has much larger mean disagreement: $17.8546$\% and
$9.3960$\%. Its target-one score-RMS/pulse ratio has mean approximately
$911.2$ but median only $1.065$, indicating strong sensitivity to exceptional
seeds. Treating this as typical persistent personalization damage would
ignore both unstable scale and a deeper issue: the meaning of ``the same
noise'' changes when factor coordinates are rebalanced.

\subsection{Orthogonal alignment control}
Equivalent factors can be rotated as $(A,B)\mapsto(AR,R^TB)$ for orthogonal
$R$ without changing $Q$. E2b aligns a branch to the current reference before
each continuation round. It solves the orthogonal Procrustes problem using
the SVD of
\begin{equation}
 A_{\mathrm{branch}}^TA_{\mathrm{ref}}+
 B_{\mathrm{branch}}B_{\mathrm{ref}}^T=U\Sigma V^T,
 \qquad R=UV^T.
 \label{eq:alignment}
\end{equation}
The transformed branch uses the same future coordinate draws in this aligned
basis. Orthogonal transformations preserve isotropic Gaussian marginal laws
in exact arithmetic, but alter the joint coupling between branches.
Consequently a smaller paired distance does not imply an improved marginal
training algorithm. Alignment here needs a counterfactual reference and is
an oracle control, not a deployable intervention.

E2b completes 40 declared Two-r8 probes. Maximum measured relative product
error from alignment is $3.9133\times10^{-8}$. At full amplitude, the lag-ten
shared-reset disagreement falls to $0.4204$\% and $0.1083$\% at targets
one and two. The large raw-coordinate distances are therefore highly
coupling-sensitive. This does not identify alignment as a uniquely correct
causal coupling: it demonstrates that the original interpretation was not
invariant to a product-preserving representation choice.

\reportfigure{05_local_state_memory.pdf}{Pulse/reset trajectories}
{E2/E2b top-10 disagreement for retained pulse, shared reset, and local reset
at starting target epsilon one and pulse amplitude one. Means and SDs use
five checkpoint seeds after averaging two draws within seed. Candidate sets
exclude training positives. The interventions access private states and
measure prediction disagreement, not relevance or accuracy gains.}{fig:memory}
\reportfigure{06_coupling_control.pdf}{The factor-coordinate coupling control}
{Two-r8 shared-reset disagreement at lag ten under raw and aligned coupling.
Diamonds show means, whiskers seed SD, and circles all five seed values.
The logarithmic axis reports percent top-10 set disagreement. Differences
of 17.8546 versus 0.4204 percent at target one, and 9.3960 versus 0.1083
percent at target two, concern joint diagnostic coupling.}{fig:coupling}

\subsection{Verification and limits}
Six E2 and two E2b synthetic cases check zero-pulse equality, both-reset
equality, RNG matching, reset identities, score closure, and product
preservation. In real probes, both-reset states are constructed and checked
at reset; the ten-round both-reset equality test belongs to the synthetic
checks. This distinction avoids overstating the real experiment's controls.

The aligned results leave modest local-state effects and an important
measurement concern. E2b is an adaptive follow-up prompted by E2, disclosed
as such. No held-out relevance evaluation or new privacy claim accompanies
it. Continued releases would need additional composition beyond the starting
epsilon. E3 tests whether the coupling phenomenon persists on a larger
dataset and whether the balancing operation explains its size.
\FloatBarrier
